% GENERATED FILE. Edit canonical lessons, then run npm run generate:books.
% canonical-source-hash: fnv1a64:48344604a6ce3cbf
% canonical-lessons: TE-S01-letter-ta, TE-C05-maatlaadu, TE-S147-digit-8, TE-C05-nenu-telugu-maatlaadataanu, TE-S111-letter-sa, TE-C05-undu, TE-S125-letter-ha, TE-C05-pani-ceyu, TE-S148-digit-9, TE-C05-practice

\chapter{The First Verbs}
\label{ch:first-verbs}
\hlchaptermodality{5}

\begin{chapteropening}
\textbf{By the end of this chapter:} \emph{I can build short present-tense Telugu sentences with the chapter's first verbs.}

Until now, “to be” and “to be-not.” Now verbs that act, and Telugu's sentences begin to move: \te{నేను} \te{తెలుగు} \te{మాట్లాడతాను} — stem + tense + person, stacked.
\end{chapteropening}

\section[ta]{\te{త} — one character, met inside words you already say}

\label{lesson:TE-S01-letter-ta}



\begin{quote}
Before the new one: \te{౭} — what amount does it stand for? And one from further back: \te{డ}?

One character this time — and you have been saying it for pages without knowing which mark on the page it was.
\end{quote}

\subsection*{Script you'll notice: \te{త}}
\textbf{\te{త}} — \emph{ta}.

It is a \textbf{consonant}, and in this script a consonant is never bare: it comes with an \emph{a} already in it. So it is not \emph{t}, it is \textbf{ta}.

You already say these, and every one of them has \te{త} somewhere inside it:

\begin{itemize}
\raggedright
  \item \textbf{\te{సంతోషం}} — joy / pleased to meet you
  \item \textbf{\te{వెళ్ళి} \te{వస్తాను}} \emph{veḷḷi vastānu} — goodbye (lit. \textquotedblleft{}I'll go and come back\textquotedblright{})
\end{itemize}

\subsection*{Writing: \te{త} — copy what you see}
Put your pen on \te{త} and follow its line. Copy the shape you can see — slowly, and larger than it is printed.

\begin{quote}
This book does not yet tell you \textbf{where to start the character or which way to travel}. That is a real thing, taught with real variation from school to school, and it is not written down here until it can be written down with a source. Copying what is in front of you needs no such source, and it is how the shape gets into your hand in the meantime.
\end{quote}

\subsection*{Your turn}
\begin{itemize}
\raggedright
  \item \emph{Look:} at these words, and find \te{త} in the ones that have it
\end{itemize}

\begin{quote}
\te{సంతోషం}  ·  \te{వెళ్ళి} \te{వస్తాను}  ·  \te{నమస్కారం}
\end{quote}

\begin{itemize}
\raggedright
  \item \emph{Trace it:} \te{త} three times, saying \emph{ta} as you finish each one
  \item \emph{Look:} back at any page of this chapter and find \te{త} once more
\end{itemize}

\begin{tcolorbox}[breakable,colback=teal!4,colframe=teal!35!black,title={Before you move on}]
Which character is this — \te{త}? What sound does it carry? (\emph{\textbf{ta}}.) Name one word you already say that contains it.

One more, from much earlier: \te{౨} — what amount does it stand for?
\end{tcolorbox}
\section[māṭlāḍu]{\te{మాట్లాడు} (māṭlāḍu) — \textquotedblleft{}to speak,\textquotedblright{} your first full verb}

\label{lesson:TE-C05-maatlaadu}



\begin{quote}
So far, \textquotedblleft{}to be.\textquotedblright{} Here is a verb that \emph{does} something — and it shows how Telugu builds \textquotedblleft{}I do X.\textquotedblright{}
\end{quote}

\subsection*{The letters in this word}
\textbf{\te{మా}} (\emph{mā}) + \textbf{\te{ట్లా}} (\emph{ṭlā}) + \textbf{\te{డు}} (\emph{ḍu}) $\to$ \textbf{\te{మాట్లాడు}} (\emph{māṭlāḍu}).

\begin{cousinweb}
\textbf{\te{మాట్లాడు}} (\emph{māṭlāḍu}, \textquotedblleft{}to speak, to talk\textquotedblright{}) is native Dravidian, built on \textbf{\te{మాట}} (\emph{māṭa}, \textquotedblleft{}word, speech\textquotedblright{}). To say \textquotedblleft{}\textbf{I speak},\textquotedblright{} Telugu adds a two-part ending to the stem — tense + person:

\begin{quote}
\emph{māṭlāḍ-} (speak) + \emph{-t-} (present/future) + \emph{-ānu} (\textquotedblleft{}I\textquotedblright{}) $\to$ \textbf{\te{మాట్లాడతాను}} (\emph{māṭlāḍatānu}), \textquotedblleft{}I speak / I will speak.\textquotedblright{}
\end{quote}

The \textbf{-\te{ాను}} (\emph{-ānu}) ending \emph{is} \textquotedblleft{}I,\textquotedblright{} so \emph{nēnu} is rarely needed. Change the ending, change the person: \emph{māṭlāḍtāḍu} (he), \emph{māṭlāḍtundi} (she/it), \emph{māṭlāḍtāru} (they / you-respectful).
\end{cousinweb}

\begin{grammarlens}[title={Grammar Lens: stem + tense + person}]
Every Telugu verb is built this way — a stem, a tense-marker, a person-ending, stacked and glued. Learn the pattern once and every verb follows it. (Telugu, like Tamil, marks no gender in the first person: \emph{māṭlāḍatānu} is the same for a man or a woman.)
\end{grammarlens}

\begin{sounds}
\te{మా} (\emph{mā}) + \te{ట్లా} (\emph{ṭlā}) + \te{డు} (\emph{ḍu}) $\to$ \te{మాట్లాడు}.
\end{sounds}

\subsection*{Your turn}
Say these aloud:

\begin{itemize}
\raggedright
  \item \textquotedblleft{}māṭlāḍu\textquotedblright{} (speak)
  \item \textquotedblleft{}I speak\textquotedblright{} — \emph{māṭlāḍatānu}
  \item the word \textquotedblleft{}speech\textquotedblright{} inside it (\emph{māṭa})
\end{itemize}

\begin{tcolorbox}[breakable,colback=teal!4,colframe=teal!35!black,title={Before you move on}]
Build \textquotedblleft{}I speak\textquotedblright{} from \emph{māṭlāḍu}, and name the noun it's built on. (\emph{Māṭlāḍatānu}; \emph{māṭa}, \textquotedblleft{}word / speech.\textquotedblright{})
\end{tcolorbox}
\section[8]{\te{౮} — the Telugu numeral for 8}

\label{lesson:TE-S147-digit-8}



\begin{quote}
Before the new one: \te{త} — what does it do? And one from further back: \te{౬}?

One numeral this time, and it is read the way every numeral is read: as an amount.
\end{quote}

\subsection*{Script you'll notice: \te{౮}}
\textbf{\te{౮}} is \textbf{eight}.

Eight of ten. Read the row below and say each value aloud before moving on — the pause matters more than the speed.

\subsection*{Writing: \te{౮} — copy what you see}
Put your pen on \te{౮} and follow its line. Copy the shape you can see — slowly, and larger than it is printed.

\begin{quote}
This book does not yet tell you \textbf{where to start the character or which way to travel}. That is a real thing, taught with real variation from school to school, and it is not written down here until it can be written down with a source. Copying what is in front of you needs no such source, and it is how the shape gets into your hand in the meantime.
\end{quote}

\subsection*{Your turn}
\begin{itemize}
\raggedright
  \item \emph{Look:} at this row and say the value of each numeral in it, left to right
\end{itemize}

\begin{quote}
\te{౮}  ·  \te{౩}  ·  \te{౭}  ·  \te{౧}
\end{quote}

\begin{itemize}
\raggedright
  \item \emph{Trace it:} \te{౮} three times, saying \textquotedblleft{}8\textquotedblright{} as you finish each one
  \item \emph{Look:} for \te{౮} on any page of a Telugu book, calendar or doorplate you can find
\end{itemize}

\begin{tcolorbox}[breakable,colback=teal!4,colframe=teal!35!black,title={Before you move on}]
Which numeral is this — \te{౮}? What amount does it stand for? (\te{౮} is 8.)

One more, from much earlier: \te{◌ి} — what does it do?
\end{tcolorbox}
\section[nēnu telugu māṭlāḍatānu]{\te{నేను} \te{తెలుగు} \te{మాట్లాడతాను} (nēnu telugu māṭlāḍatānu) — \textquotedblleft{}I speak Telugu\textquotedblright{}}

\label{lesson:TE-C05-nenu-telugu-maatlaadataanu}



\begin{quote}
Your first full, moving sentence — and it names one of India's most musical languages.
\end{quote}

\subsection*{The letters in this word}
\textbf{\te{తెలుగు}} (\emph{telugu}): \textbf{\te{తె}} (\emph{te}) + \textbf{\te{లు}} (\emph{lu}) + \textbf{\te{గు}} (\emph{gu}).

\begin{cousinweb}
\textbf{\te{నేను} \te{తెలుగు} \te{మాట్లాడతాను}} = \textbf{\te{నేను}} (\emph{nēnu}, \textquotedblleft{}I\textquotedblright{}) + \textbf{\te{తెలుగు}} (\emph{telugu}, \textquotedblleft{}Telugu,\textquotedblright{} the object) + \textbf{\te{మాట్లాడతాను}} (\emph{māṭlāḍatānu}, \textquotedblleft{}speak\textquotedblright{}) — \textquotedblleft{}\textbf{I Telugu speak},\textquotedblright{} the verb last as always. The name \textbf{\te{తెలుగు}} (\emph{telugu}) is native, of debated origin (often linked to \emph{tenugu}, and to the land \emph{Trilinga}); its vowel-rich sound and habit of ending most words in a vowel earned it the old nickname \textquotedblleft{}the Italian of the East.\textquotedblright{} It is the language of the \emph{Kāvya} poets and, today, of some eighty million speakers.
\end{cousinweb}

\begin{grammarlens}[title={Grammar Lens: no gender on \textquotedblleft{}I speak\textquotedblright{}}]
\emph{Māṭlāḍatānu} is the same whether a man or a woman says it — Telugu, like Tamil, marks \textbf{no gender in the first or second person} (only the third: \emph{-āḍu} he, \emph{-undi} she/it). This is unlike Hindi, where even \textquotedblleft{}I speak\textquotedblright{} changes for the speaker's gender (\emph{boltā/boltī}).
\end{grammarlens}

\subsection*{Your turn}
Say these aloud:

\begin{itemize}
\raggedright
  \item \textquotedblleft{}nēnu telugu māṭlāḍatānu\textquotedblright{}
  \item the old nickname for Telugu and why (\textquotedblleft{}the Italian of the East\textquotedblright{} — vowel-rich, most words end in a vowel)
  \item does \textquotedblleft{}I speak\textquotedblright{} change for a man vs. a woman? (No — no 1st-person gender)
\end{itemize}

\begin{tcolorbox}[breakable,colback=teal!4,colframe=teal!35!black,title={Before you move on}]
Say \textquotedblleft{}I speak Telugu,\textquotedblright{} and give Telugu's old nickname. (\emph{Nēnu telugu māṭlāḍatānu}; \textquotedblleft{}the Italian of the East.\textquotedblright{})
\end{tcolorbox}
\section[sa]{\te{స} — one character, met inside words you already say}

\label{lesson:TE-S111-letter-sa}



\begin{quote}
Before the new one: \te{౮} — what amount does it stand for? And one from further back: \te{ప}?

One character this time — and you have been saying it for pages without knowing which mark on the page it was.
\end{quote}

\subsection*{Script you'll notice: \te{స}}
\textbf{\te{స}} — \emph{sa}.

It is a \textbf{consonant}, and in this script a consonant is never bare: it comes with an \emph{a} already in it. So it is not \emph{s}, it is \textbf{sa}.

You already say these, and every one of them has \te{స} somewhere inside it:

\begin{itemize}
\raggedright
  \item \textbf{\te{నమస్కారం}} \emph{namaskāram} — hello / greetings (namaskāram — \textquotedblleft{}a making of a bow\textquotedblright{})
  \item \textbf{\te{సరే}} \emph{sarē} — okay / alright (sarē)
  \item \textbf{\te{వెళ్ళి} \te{వస్తాను}} \emph{veḷḷi vastānu} — goodbye (lit. \textquotedblleft{}I'll go and come back\textquotedblright{})
\end{itemize}

\subsection*{Writing: \te{స} — copy what you see}
Put your pen on \te{స} and follow its line. Copy the shape you can see — slowly, and larger than it is printed.

\begin{quote}
This book does not yet tell you \textbf{where to start the character or which way to travel}. That is a real thing, taught with real variation from school to school, and it is not written down here until it can be written down with a source. Copying what is in front of you needs no such source, and it is how the shape gets into your hand in the meantime.
\end{quote}

\subsection*{Your turn}
\begin{itemize}
\raggedright
  \item \emph{Look:} at these words, and find \te{స} in the ones that have it
\end{itemize}

\begin{quote}
\te{నమస్కారం}  ·  \te{సరే}  ·  \te{ధన్యవాదములు}
\end{quote}

\begin{itemize}
\raggedright
  \item \emph{Trace it:} \te{స} three times, saying \emph{sa} as you finish each one
  \item \emph{Look:} back at any page of this chapter and find \te{స} once more
\end{itemize}

\begin{tcolorbox}[breakable,colback=teal!4,colframe=teal!35!black,title={Before you move on}]
Which character is this — \te{స}? What sound does it carry? (\emph{\textbf{sa}}.) Name one word you already say that contains it.

One more, from much earlier: \te{◌ం} — what does it do?
\end{tcolorbox}
\section[uṇḍu]{\te{ఉండు} (uṇḍu) — \textquotedblleft{}to be, to stay, to live\textquotedblright{}}

\label{lesson:TE-C05-undu}



\begin{quote}
You have already used this verb — in \textquotedblleft{}how are you.\textquotedblright{} Now meet it in its own right, and use it to say where you live.
\end{quote}

\subsection*{The letters in this word}
\textbf{\te{ఉ}} (independent \emph{u}) + \textbf{\te{ండు}} (\emph{ṇḍu}, with the retroflex cluster \textbf{\te{ండ}}) $\to$ \textbf{\te{ఉండు}} (\emph{uṇḍu}).

\begin{cousinweb}
\textbf{\te{ఉండు}} (\emph{uṇḍu}, \textquotedblleft{}to be, to exist, to stay, to live\textquotedblright{}) is native Dravidian and does an enormous amount of work. It is the verb behind Chapter 3's \emph{unnāru} (\textquotedblleft{}you are\textquotedblright{}) and \emph{bāgunnānu} (\textquotedblleft{}I'm well\textquotedblright{}) — and it also means \textquotedblleft{}to reside\textquotedblright{}: \textbf{\te{నేను} \te{హైదరాబాద్లో} \te{ఉంటాను}} (\emph{nēnu Hyderābādlō uṇṭānu}, \textquotedblleft{}I live in Hyderabad\textquotedblright{}). One verb covers \textquotedblleft{}be,\textquotedblright{} \textquotedblleft{}stay,\textquotedblright{} and \textquotedblleft{}live.\textquotedblright{}
\end{cousinweb}

\begin{grammarlens}[title={Grammar Lens: \textquotedblleft{}in\textquotedblright{} comes after the noun}]
\textbf{\te{హైదరాబాద్లో}} (\emph{Hyderābādlō}) = \emph{Hyderabad} + \textbf{-\te{లో}} (\emph{-lō}, \textquotedblleft{}in\textquotedblright{}). The little word for \textquotedblleft{}in\textquotedblright{} is glued to the \textbf{end} of the noun — Telugu uses \textbf{post}positions (really, case-suffixes) where English uses prepositions. \textquotedblleft{}Hyderabad-in I stay,\textquotedblright{} the verb last.
\end{grammarlens}

\subsection*{Your turn}
Say these aloud:

\begin{itemize}
\raggedright
  \item \textquotedblleft{}uṇḍu\textquotedblright{}
  \item \textquotedblleft{}I live in Hyderabad\textquotedblright{} — \emph{nēnu Hyderābādlō uṇṭānu}
  \item where \textquotedblleft{}in\textquotedblright{} (\emph{-lō}) sits — glued to the end of the noun
\end{itemize}

\begin{tcolorbox}[breakable,colback=teal!4,colframe=teal!35!black,title={Before you move on}]
Name three things the one verb \emph{uṇḍu} can mean, and say where \textquotedblleft{}in\textquotedblright{} goes. (\textquotedblleft{}To be / to stay / to live\textquotedblright{}; \emph{-lō} \textquotedblleft{}in\textquotedblright{} attaches to the \textbf{end} of the noun.)
\end{tcolorbox}
\section[ha]{\te{హ} — one character, met inside words you already say}

\label{lesson:TE-S125-letter-ha}



\begin{quote}
Before the new one: \te{స} — what does it do? And one from further back: \te{౭}?

One character this time — and you have been saying it for pages without knowing which mark on the page it was.
\end{quote}

\subsection*{Script you'll notice: \te{హ}}
\textbf{\te{హ}} — \emph{ha}.

It is a \textbf{consonant}, and in this script a consonant is never bare: it comes with an \emph{a} already in it. So it is not \emph{h}, it is \textbf{ha}.

You met it on the page before this one, at the front of a place name:

\begin{itemize}
\raggedright
  \item \textbf{\te{హైదరాబాద్}} \emph{Hyderābād} — the city from the sentence about where you live
\end{itemize}

Look at where that word came from. \emph{Hyderābād} reached Telugu through Persian and Urdu, and that is the pattern: in the words this book has given you so far, \te{హ} sits almost entirely on borrowings — from Persian, from Urdu, from Sanskrit. It is a letter you will meet again on the very next page, in a word for help that came in from Sanskrit.

\subsection*{Writing: \te{హ} — copy what you see}
Put your pen on \te{హ} and follow its line. Copy the shape you can see — slowly, and larger than it is printed.

\begin{quote}
This book does not yet tell you \textbf{where to start the character or which way to travel}. That is a real thing, taught with real variation from school to school, and it is not written down here until it can be written down with a source. Copying what is in front of you needs no such source, and it is how the shape gets into your hand in the meantime.
\end{quote}

\subsection*{Your turn}
\begin{itemize}
\raggedright
  \item \emph{Look:} at these words, and find \te{హ} in the ones that have it
\end{itemize}

\begin{quote}
\te{హైదరాబాద్}  ·  \te{నమస్కారం}
\end{quote}

\begin{itemize}
\raggedright
  \item \emph{Trace it:} \te{హ} three times, saying \emph{ha} as you finish each one
  \item \emph{Look:} back at any page of this chapter and find \te{హ} once more
\end{itemize}

\begin{tcolorbox}[breakable,colback=teal!4,colframe=teal!35!black,title={Before you move on}]
Which character is this — \te{హ}? What sound does it carry? (\emph{\textbf{ha}}.) Name one word you already say that contains it.

One more, from much earlier: \te{౩} — what amount does it stand for?
\end{tcolorbox}
\section[pani cēyu]{\te{పని} \te{చేయు} (pani cēyu) — \textquotedblleft{}to work\textquotedblright{}}

\label{lesson:TE-C05-pani-ceyu}



\begin{quote}
The last verb of the chapter — and, like Hindi's \emph{karnā} and Tamil's \emph{sey}, it is a \textquotedblleft{}do\textquotedblright{} verb that builds a hundred others.
\end{quote}

\subsection*{The letters in this word}
\textbf{\te{పని}} (\emph{pani}, \textquotedblleft{}work\textquotedblright{}): \textbf{\te{ప}} (\emph{pa}) + \textbf{\te{ని}} (\emph{ni}). \textbf{\te{చేయు}} (\emph{cēyu}, \textquotedblleft{}do\textquotedblright{}): \textbf{\te{చే}} (\emph{cē}) + \textbf{\te{యు}} (\emph{yu}).

\begin{cousinweb}
\textbf{\te{చేయు}} (\emph{cēyu}, \textquotedblleft{}to do, to make\textquotedblright{}) is a core native Dravidian verb, and Telugu — like Hindi with \emph{karnā} and Tamil with \emph{sey} — builds compound verbs by putting a noun in front of it:

\begin{itemize}
\raggedright
  \item \textbf{\te{పని} \te{చేయు}} (\emph{pani cēyu}) = \textquotedblleft{}work-do\textquotedblright{} = \textquotedblleft{}\textbf{to work}\textquotedblright{} (\emph{pani}, \textquotedblleft{}work\textquotedblright{})
  \item \textbf{\te{వంట} \te{చేయు}} (\emph{vaṇṭa cēyu}) = \textquotedblleft{}to cook\textquotedblright{}
  \item \textbf{\te{సహాయం} \te{చేయు}} (\emph{sahāyaṁ cēyu}) = \textquotedblleft{}to help\textquotedblright{} (\emph{sahāyaṁ} $\leftarrow$ Sanskrit)
\end{itemize}

So \textquotedblleft{}I work\textquotedblright{} is \textbf{\te{నేను} \te{పని} \te{చేస్తాను}} (\emph{nēnu pani cēstānu}) — \emph{cēy-} + tense + \textquotedblleft{}I.\textquotedblright{} Master \emph{cēyu} and a whole grammar of \textquotedblleft{}do-verbs\textquotedblright{} opens up — and notice \emph{sahāyaṁ cēyu} (\textquotedblleft{}to help\textquotedblright{}) pairs a \textbf{Sanskrit} noun with the \textbf{native} \emph{cēyu}: Telugu's two vocabularies working together in a single verb.
\end{cousinweb}

\begin{grammarlens}[title={Grammar Lens: noun + \te{చేయు} = a new verb}]
This \textquotedblleft{}noun + \emph{cēyu}\textquotedblright{} pattern is the exact twin of Hindi's \textquotedblleft{}noun + \emph{karnā}\textquotedblright{} and Tamil's \textquotedblleft{}noun + \emph{sey}.\textquotedblright{} Three languages — two of them from unrelated families — all hit on the same trick: an all-purpose \textquotedblleft{}do\textquotedblright{} verb with nouns bolted on.
\end{grammarlens}

\subsection*{Your turn}
Say these aloud:

\begin{itemize}
\raggedright
  \item \textquotedblleft{}pani cēyu\textquotedblright{} (to work)
  \item \textquotedblleft{}I work\textquotedblright{} — \emph{nēnu pani cēstānu}
  \item the Hindi and Tamil twins of this trick (\emph{karnā}, \emph{sey})
\end{itemize}

\begin{tcolorbox}[breakable,colback=teal!4,colframe=teal!35!black,title={Before you move on}]
How does Telugu turn \textquotedblleft{}work\textquotedblright{} into \textquotedblleft{}to work,\textquotedblright{} and what verbs work the same way in Hindi and Tamil? (\emph{Pani} + \emph{cēyu} = \textquotedblleft{}work-do\textquotedblright{}; Hindi \emph{karnā}, Tamil \emph{sey}.)
\end{tcolorbox}
\section[9]{\te{౯} — the Telugu numeral for 9}

\label{lesson:TE-S148-digit-9}



\begin{quote}
Before the new one: \te{హ} — what does it do? And one from further back: \te{త}?

One numeral this time, and it is read the way every numeral is read: as an amount.
\end{quote}

\subsection*{Script you'll notice: \te{౯}}
\textbf{\te{౯}} is \textbf{nine}.

That is nine of the ten, and the last one is the one that makes the rest useful.

\subsection*{Writing: \te{౯} — copy what you see}
Put your pen on \te{౯} and follow its line. Copy the shape you can see — slowly, and larger than it is printed.

\begin{quote}
This book does not yet tell you \textbf{where to start the character or which way to travel}. That is a real thing, taught with real variation from school to school, and it is not written down here until it can be written down with a source. Copying what is in front of you needs no such source, and it is how the shape gets into your hand in the meantime.
\end{quote}

\subsection*{Your turn}
\begin{itemize}
\raggedright
  \item \emph{Look:} at this row and say the value of each numeral in it, left to right
\end{itemize}

\begin{quote}
\te{౯}  ·  \te{౮}  ·  \te{౬}  ·  \te{౨}  ·  \te{౫}
\end{quote}

\begin{itemize}
\raggedright
  \item \emph{Trace it:} \te{౯} three times, saying \textquotedblleft{}9\textquotedblright{} as you finish each one
  \item \emph{Look:} for \te{౯} on any page of a Telugu book, calendar or doorplate you can find
\end{itemize}

\begin{tcolorbox}[breakable,colback=teal!4,colframe=teal!35!black,title={Before you move on}]
Which numeral is this — \te{౯}? What amount does it stand for? (\te{౯} is 9.)

One more, from much earlier: \te{ఉ} — what does it do?
\end{tcolorbox}
\section[Practice]{Chapter 5 — The first verbs}

\label{lesson:TE-C05-practice}



\begin{quote}
No new words. Three verbs, one engine — build real sentences about yourself.
\end{quote}

\subsection*{Your turn: build the sentences}
Three sentences about yourself, read off the page before you say them:

\noindent
\begin{tabularx}{\linewidth}{@{}>{\raggedright\arraybackslash}X>{\raggedright\arraybackslash}X>{\raggedright\arraybackslash}X@{}}
\toprule
\textbf{Telugu} & \textbf{} & \textbf{English} \\
\midrule
\textbf{\te{నేను} \te{తెలుగు} \te{మాట్లాడతాను}.} & \emph{nēnu telugu māṭlāḍatānu} & I speak Telugu. \\
\textbf{\te{నేను} \te{హైదరాబాద్లో} \te{ఉంటాను}.} & \emph{nēnu Hyderābādlō uṇṭānu} & I live in Hyderabad. \\
\textbf{\te{నేను} \te{పని} \te{చేస్తాను}.} & \emph{nēnu pani cēstānu} & I work. \\
\bottomrule
\end{tabularx}

A few of these characters have not had their own lesson yet. They are here to be \textbf{recognised}, not decoded — the romanization beside each line is what you read from. The letters catch up in the chapters ahead.

Say these aloud:

\begin{itemize}
\raggedright
  \item \textquotedblleft{}I speak Telugu\textquotedblright{} — \emph{nēnu telugu māṭlāḍatānu}
  \item \textquotedblleft{}I live in Hyderabad\textquotedblright{} — \emph{nēnu Hyderābādlō uṇṭānu}
  \item \textquotedblleft{}I work\textquotedblright{} — \emph{nēnu pani cēstānu}
  \item \textquotedblleft{}he speaks\textquotedblright{} vs. \textquotedblleft{}she speaks\textquotedblright{} — \emph{māṭlāḍtāḍu} / \emph{māṭlāḍtundi}
\end{itemize}

\begin{grammarlens}[title={What you've built — the engine}]
Say these aloud:

\begin{itemize}
\raggedright
  \item the three stacked pieces of every Telugu verb (stem + tense + person)
  \item the ending that means \textquotedblleft{}I\textquotedblright{} (\emph{-ānu})
  \item where Telugu marks gender — and where it doesn't (3rd person only)
\end{itemize}
\end{grammarlens}

\begin{cousinweb}
\textbf{Roots you now carry:}

\begin{itemize}
\raggedright
  \item \emph{māṭlāḍu} \textquotedblleft{}to speak\textquotedblright{} $\leftarrow$ \emph{māṭa} \textquotedblleft{}word\textquotedblright{}
  \item \emph{uṇḍu} \textquotedblleft{}to be / stay / live\textquotedblright{}; \emph{-lō} \textquotedblleft{}in\textquotedblright{} attaches to the noun's end
  \item \emph{cēyu} \textquotedblleft{}to do\textquotedblright{} $\to$ \emph{pani cēyu} \textquotedblleft{}to work\textquotedblright{} (the twin of Hindi's \emph{karnā}, Tamil's \emph{sey})
  \item \emph{telugu} — \textquotedblleft{}the Italian of the East,\textquotedblright{} vowel-rich
\end{itemize}
\end{cousinweb}

\begin{tcolorbox}[breakable,colback=teal!4,colframe=teal!35!black,title={Before you move on}]
In one breath, say your language, your city, and your work in Telugu. (\emph{Nēnu telugu māṭlāḍatānu. Nēnu … lō uṇṭānu. Nēnu pani cēstānu.})
\end{tcolorbox}
